1. The Lotus-Like Lay Buddhist

A person who has gone to the Buddha for refuge and who sincerely wishes to become a good lay disciple should follow the advice taught by the Enlightened One himself in the Discourses that he delivered to his disciples. Although the Buddha had passed away over 2500 years ago, fortunately, these Discourses (Sutta) have been preserved intact by his noble disciples in the Buddhist Scriptures called the Tipitaka or Three Baskets.

Today many people prefer to go into the Internet to access their information. If one tries to look for the answer there, one will surely find all sorts of suggestions on what makes a good ‘Buddhist’. Many of these suggestions are based on the personal ideas of various writers, interpretations of popular culture or idealized notions of what a ‘good Buddhist’ should be. Secondly, the so-called ‘marks’ characterising a good Buddhist varies according to the inclination of the writers, and can lead to confusion. The objective of this book is to explain to the ordinary layperson what the Buddha taught about the qualities that make a Lotus-like Lay Buddhist a good Buddhist. In the Candala Sutta (Discourse on Outcaste) in Anguttara Nikaya V, the Buddha said:

“Bikkhus, pursuing five things, a layman is an outcaste of laymen, a stain of a laymen, a dreg of laymen. Which five? (1) He/she does not have confidence (in the Buddha); (2) is not virtuous; (3) is eager for protective charms & ceremonies; (4) trusts protective charms and ceremonies, not kamma; (5) seeks outside (the Sangha) for a gift-worthy person, and offers there first.”

1.1 Five Qualities of the Lotus-like Lay Buddhist

Alternately, in the same Discourse, the Buddha lists five things that make a Lotus-like layperson: “Bikkhus (see 1.3), pursuing five things, a layman is a jewel of laymen, a lotus of laymen, a fine flower of laymen. Which five?
(1) He/she had confidence (in the Buddha).
(2) He/She is virtuous.
(3) He/She is not eager for protective charms and ceremonies (not superstitious)
(4) He/She trusts kamma, not protective charms and ceremonies
(5) He/She seeks within (the Sangha) for a gift-worthy person, and offers here first.”

1.2 Relationship between the Sangha and the Laity

Although this discourse is about the qualities of laymen, yet the Buddha taught it to the monks so that they can preach it to benefit the laity. The relationship between the laity and members of the Sangha (monkhood) is one of mutual benefit. The material needs of the monk is provided by the laity while the monkhood look after the spiritual needs of the laity. This unique relationship is well defined by the Buddha in his Discourse to the youth Sigala as follows:

“In five ways, young householder, should a householder minister to the samana-brahmana or ascetics and Brahmins (see 1.3) as the Zenith (i) by lovable deeds, (ii) by lovable words, (iii) by lovable thoughts, (iv) by keeping open house to them, (v) by supplying their material needs.

“The samana-brahmana thus ministered to as the Zenith show their compassion toward him in six ways: (i) they restrain him from evil, (ii) they persuade him to do good, (iii) they love him with a kind heart, (iv) they make him hear what he has not heard, (v) they clarify what he has already heard, (vi) they point out the path to a heavenly state.

“In these six ways do the samana-brahmana show their compassion torwards a householder who ministers to them as the Zenith. Thus is the Zenith covered by him and made safe and secure.” (Sigalovada Sutta, Digha Nakaya).

1.3 Meanings of the terms ‘bhikkhu’ and ‘samana-brahmana’

1. The Buddha used the term ‘bhikkhu’, which means ‘almsman or mendicant’ to address his male disciples who had renounced the household life to join his Order of monks or Bhikkhu Sangha. During the Buddha’s time, bhikkhus were recognized by others by the token of their begging bowls or ‘patta’, which confers sanctity as the visible symbol and token of the almsman’s calling. People were aware of their difference in standing from vagabonds and beggars and recognized that they were in quest of spiritual liberation. For his daily subsistence, the bhikkhu had to go on almsround to receive food from the public. The food received in the almsbowl is called ‘pindapata’. He is not allowed to ask but must stand silent at the door of a house till food is offered to him. If not offering is forthcoming, he can move to another doorway and wait until he obtains pindapata. All this is done in silence with dignity and utmost restraint of the sense faculties.

The monastic code forbids the bhikkhu from receiving god, silver, or cash. The female counterpart is called ‘bhikkhuni’ and the Order of nuns is called Bhikkhuni Sangha. According to the Theravada tradition, the Bhikkhuni Sangha is extinct today.

2. The term ‘samana-brahmana’ is generally used to denote the two religious systems during the Buddha’s time. The Brahmanas adhered to the age old Vedic religion and considered their Scriptures or Vedas to by the ultimate spiritual authority. The priests, who were the custodians of such prayers, assumed a very high degree of spiritual supremacy in Vedic society by claiming kinship with their creator Brahma, thereby creating a caste system to perpetuate their status quo. Indian society under the Vedic system was based on the caste system. Next to the Brahmans were the warrior caste called Kittiyas, followed by the Vessas or merchant caste. The Sudras, the lowest caste, were considered ineligible to perform spiritual rites.

The other group of religious practitioners were the Samanas, who rejected the Vedas and most Brahmanical beliefs and practices. They were mostly wandering almsmen who usually lived in forests, striving for spiritual attainment. The term ‘Samana’ means ‘labourer in spiritual life’. In Buddhist literature the term Samana is applicable to the order of six Titthiyas, founders of heretical schools that were popular during the Buddha’s time. Although the Buddha addressed his monks as Bikkhus, the people called them Sakyaputtiya Samanas, with the Buddha himself addressed as the Great Samana.

The term ‘samana-brahmana’ is often translated as ‘ascetics and Brahmins’. However when the Buddha used it to describe his disciples, it would mean those who have attained sainthood. In Dhammapada Verse 142, the Buddha said that one is called a ‘brahmana’ or ‘samana’ because he has subdued all evil. The Buddha also used the term ‘brahmana’ to denot the Arahant, the highest level of attainment in his Dispensation just as Brahmana is the highest caste in Vedic culture. In the Sigalovada Sutta, Sigala’s parents were Brahmins but devout Buddhists both having attained Stream-Entry (Sotapanna). Before the old Brahman died, he advised his son, who had no interest in religion to worship the six quarters, knowing that the Buddha or one of his disciples would see the act and explain to his son the actual significance of his action. Thus in the Sigalovada Sutta, the term ‘samana-brahmana’ would mean the Ariya or Noble Ones.

In the Samana Brahmana Sutta of the Nidana Samyutta, Dasabala Vagg, the Buddha declared that only those samanas and brahmanas who comprehend the Four Noble Truths are deserving of being called true samanas and brahmanas!
